\section{Polynomial descriptions characterize polynomial advice}
The class \(\Ppoly\) consists of languages decidable in polynomial time
with advice of polynomial length depending only on input length.
Equivalently, these are languages with polynomial-size Boolean circuit
families: hardwire advice into computation circuits, or evaluate an
encoded circuit supplied as advice \cite{AroraBarak2009}.

\begin{theorem}[Advice characterization \claimid{T01}]
\label{thm:advice}
Under the evaluator convention of Definition~\ref{def:local},
\[
 \SAT\in\Ppoly
 \quad\Longleftrightarrow\quad
 \exists a,b\in\mathbb N_{>0}\ \exists n_0\
 \forall n\geq n_0:\quad \loc_{n^b}(T_n)\leq n^a.
\]
\end{theorem}
\begin{proof}
Suppose first that the right-hand side holds. For each \(n\geq n_0\),
choose a witnessing description \(d_n\). Use \(d_n\) as advice and run
the fixed evaluator \(U(d_n,x)\). Both advice length and time are
polynomial in \(n\). The finitely many smaller input lengths can be
handled by a finite table in the deciding machine. Thus
\(\SAT\in\Ppoly\).

Conversely, suppose a fixed deterministic machine \(M\) decides
SAT using advice \(a_n\) of length at most \(q(n)\) in time at most
\(p(n)\), for polynomials \(p,q\). Encode \(M\), a self-delimiting
header, and \(a_n\) in \(d_n\). Its length is polynomial in \(n\).
Efficient universality supplies a polynomial time bound for
\(U(d_n,x)\). Increasing integer exponents and discarding finitely
many small \(n\) absorbs multiplicative constants, headers, and
simulation overhead.
\end{proof}

This is a polynomial-resource equivalence, not equality of program
length and gate count. The proof applies equally to any language.

\begin{corollary}[The exact lower-bound target \claimid{C01}]
\label{cor:target}
The following assertions are equivalent:
\begin{enumerate}[label=(\roman*)]
 \item \(\SAT\notin\Ppoly\).
 \item For every \(a,b\in\mathbb N_{>0}\) and every \(n_0\), there is
 \(n\geq n_0\) such that \(\loc_{n^b}(T_n)>n^a\).
\end{enumerate}
Either assertion implies \(\PP\ne\NP\).
\end{corollary}
\begin{proof}
Negating the quantified right-hand side of Theorem~\ref{thm:advice}
gives (ii). If \(\PP=\NP\), SAT has a polynomial-time algorithm and
therefore polynomial-size circuits, contradicting (i).
\end{proof}

Assertion (ii) quantifies over infinitely many input lengths, rather
than all sufficiently large lengths. An eventual lower bound would
also suffice but is not the logical negation above.
The converse implication from \(\PP\ne\NP\) to (i) is not established.
The target excludes nonuniform advice as well as uniform algorithms.

A single fixed description correctly answering all input lengths in
polynomial time is an ordinary polynomial-time SAT algorithm.
A new polynomial-size description at each length is extra freedom.
Ruling out only maximally compressed or logarithmic-length descriptions
would leave open longer polynomial descriptions with faster evaluation.
